\documentclass[twocolumn]{aastex701}
\usepackage{amsmath}
\usepackage{booktabs}

\begin{document}

\title{Title of the Paper}

\author[0000-0000-0000-0000]{First Author}
\affiliation{Department of Astronomy, University, City, Country}
\email{first@example.org}
\correspondingauthor{First Author}

\author{Second Author}
\affiliation{Observatory, City, Country}
\email{second@example.org}

\begin{abstract}
State the problem, what you did, and the main result in a few sentences.
\end{abstract}

\keywords{\uat{Cosmology}{343} --- \uat{Type Ia supernovae}{1728}}

\section{Introduction}
\label{sec:intro}
Introduce the problem and why it matters, then the contributions of this
paper. Cite in parentheses~\citep{perlmutter1999}, or by name, as \citet{planck2020} found, or several works at
once~\citep{perlmutter1999,riess1998}.

\section{Method}
\label{sec:method}
Describe what you did in enough detail to reproduce it. Numbered equations use
the \texttt{equation} environment:
\begin{equation}
  \label{eq:main}
  \mu = 5 \log_{10}\left(\frac{d_L}{10\,\mathrm{pc}}\right).
\end{equation}
Refer to it as Equation~\eqref{eq:main}.

\section{Results}
\label{sec:results}
Table~\ref{tab:results} and Figure~\ref{fig:overview} show the main results.

\begin{table}[htbp]
  \centering
  \caption{Results on the test set.}
  \label{tab:results}
  \begin{tabular}{lrr}
    \toprule
    Method & Accuracy (\%) & Time (s) \\
    \midrule
    Baseline & 81.2 & 12 \\
    Ours & \textbf{87.5} & 15 \\
    \bottomrule
  \end{tabular}
\end{table}

\begin{figure}[htbp]
  \centering
  \includegraphics[width=\columnwidth]{example-image}
  \caption{Overview of the method. Upload an image and use its file name here.}
  \label{fig:overview}
\end{figure}

\section{Discussion}
Interpret the results, compare them with earlier work and state the limitations.

\section{Conclusion}
Summarise the findings and what comes next.

\begin{acknowledgments}
Thank the people, facilities and funding that supported this work.
\end{acknowledgments}

\facilities{HST(STIS), Keck:I}

\software{astropy}

\bibliography{references}{}
\bibliographystyle{aasjournalv7}

\end{document}
